\section{Exact iPRA proof for the simultaneous two-row shortcuts}
\label{sec:shortcut-stuttering}

Let $P$ be either native PA194ORH or native PA194ORHB.  Both have initial
control~74, halt label~254, and the same rows used below.  Let $F(P)$ be the
machine obtained by replacing exactly
\[
\begin{array}{c|c|c}
\text{row}&P&F(P)\\\hline
(26,1)&1\,R\,42&1\,R\,91\\
(101,0)&0\,L\,26&0\,R\,17.
\end{array}
\]
All other 386 transition rows remain identical.  The new machines are called
PA194ORHF and PA194ORHBF, respectively.  These are compositions of supplied
physical mechanisms; no claim that the shortcut mechanisms themselves are
new is needed for the theorem.

The proof below establishes exact relative halting equivalence in iPRA.
It applies to every finitely supported initial tape in control~74, hence
in particular to the specified blank tape.  It does not depend on a bounded
sample of source executions or on the PA-calibration argument.

\subsection{A three-clause inductive tape invariant}

Write $C=(q,h,T)$ for a configuration, with $T(i)\in\{0,1\}$, and define
\[
\begin{split}
 \mathcal I(C)\quad\Longleftrightarrow\quad
 &(q=101\ \Rightarrow\ T(h-1)=0)\\
 {}\land{}&(q=11\ \Rightarrow\ [T(h-1)=0\ \lor\ T(h+1)=0])\\
 {}\land{}&(q=26\ \Rightarrow\
       [T(h+1)=0\ \land\ (T(h)=0\ \lor\ T(h+2)=0)]).
\end{split}                                                    \tag{S1}
\]
The position notation is only a convenient presentation.  In arithmetic,
a configuration is a control and two finite binary tape halves, so the
fixed-offset bit accesses in (S1) are primitive recursive functions.
The displayed predicate is quantifier-free, or equivalently a Boolean
primitive recursive test.  Control~74 satisfies it on every initial tape.

The complete incoming-edge facts in both parents are
\[
\begin{array}{c|l}
\text{target}&\text{all incoming rows, with write and movement}\\\hline
101&(96,0):0R\\
11&(14,0):0L,\ (38,0):0R,\ (39,0):0L,\ (56,0):0L\\
26&(11,0):0L,\ (101,0):0L.
\end{array}                                                    \tag{S2}
\]
These are exhaustively checked facts about the complete finite tables.
For a child $F(P)$, the first two lists are unchanged and the last becomes
just $(11,0):0L$.

We check preservation of (S1), covering all possible next controls.
An entry into 101 writes zero immediately to the left of the new head,
proving its clause.  An entry into 11 either writes zero to its new left
neighbor (the row from 38), or writes zero to its new right neighbor
(the rows from 14, 39 and 56).  Thus its disjunctive clause is preserved.
An entry into 26 from $(11,0)$ writes zero to its new right neighbor.
The old state-11 clause then says that the new current cell is zero or
its second right neighbor is zero, exactly the remaining state-26 clause.
In a parent, the additional entry from $(101,0)$ writes zero to the new
right neighbor, while the old state-101 clause says that the new current
cell is zero.  No other row enters any of these three controls.  These
arguments prove preservation for both parents and both simultaneous-patch
children.  Quantifier-free induction on time proves (S1) for their complete
runs from control~74.

In particular, reaching control~26 while reading~1 forces
\[
                  T(h+1)=T(h+2)=0.                           \tag{S3}
\]
The second zero is essential.  The immediate predecessor alone does not
justify it: it follows from the state-11 clause and the exclusion of the
state-101 predecessor when the current bit is~1.

\subsection{Two exact three-to-one identities}

The unchanged rows needed for the first shortcut are
\[
 (42,0):0R13,\qquad (13,0):0L91.
\]
If $q=26$ and $T(h)=1$, condition (S3) gives the complete parent path
\[
 (26,h,T)\longmapsto(42,h+1,T)
   \longmapsto(13,h+2,T)\longmapsto(91,h+1,T).
                                                               \tag{S4}
\]
Every write preserves the bit already present.  The child reaches precisely
the last configuration of (S4) in one transition.

For the second shortcut the unchanged rows are
\[
 (26,0):0R194,\qquad(194,0):0R17.
\]
If $q=101$ and $T(h)=0$, (S1) supplies $T(h-1)=0$, giving
\[
 (101,h,T)\longmapsto(26,h-1,T)
   \longmapsto(194,h,T)\longmapsto(17,h+1,T).
                                                               \tag{S5}
\]
Again every write preserves its bit, and the child reaches precisely the
last configuration in one transition.  Native control~194 in (S5) is a
working control; the native halt label is~254.  None of the intermediate
controls in (S4) or (S5) is a halt control.

Both identities hold for arbitrary tape contents outside the explicitly
listed cells.  They are exact equalities of complete configurations,
including state, head and all tape cells.  Applying the two patches
simultaneously does not interfere with either argument: the second patch
removes one predecessor of 26 without weakening (S1); the first patch
changes the successor of 26, not an incoming edge used in (S2).  The child
is not required to execute the skipped intermediate controls.

\subsection{A primitive recursive stuttering clock}

Use the usual semantic totalization at the halt label: after halting, the
run remains halted.  This adds no working control to the literal machine.
Let $C_n$ be the child configuration after $n$ transitions from the common
initial configuration $C_0$.  Define
\[
 \epsilon(C)=
  \begin{cases}
   1,&q=26\text{ and the scanned bit is }1,\text{ or }q=101
                         \text{ and the scanned bit is }0,\\
   0,&\text{otherwise},
  \end{cases}
\]
with $\epsilon=0$ at halt, and put
\[
 A(n)=\sum_{j<n}\epsilon(C_j),\qquad
             \Theta(n)=n+2A(n).                             \tag{S6}
\]
These are primitive recursive functions.  Equations (S4)--(S5), ordinary
row equality, and the halt totalization prove in iPRA that, whenever
$\mathcal I(C)$ holds,
\[
 \operatorname{step}_{F(P)}(C)
       =\operatorname{Run}_P(1+2\epsilon(C),C).              \tag{S7}
\]
Every local duration is either one or three.  Quantifier-free induction,
using (S1), (S7), and the elementary iterated-run composition equation,
gives the exact complete-configuration identity
\[
 \boxed{\quad
    \operatorname{Run}_P(\Theta(n),C_0)
       =\operatorname{Run}_{F(P)}(n,C_0),\qquad
       n\le\Theta(n)\le3n.\quad}                           \tag{S8}
\]
Thus the cumulative speed difference is exactly twice the number of
shortcut occurrences along the child run, not merely an unspecified
simulation overhead.

If the child halts at time $n$, the parent halts at the explicit time
$\Theta(n)$.  Conversely, if the parent is halted at time $m$, then
$\Theta(m)\ge m$ and absorption implies that the left side of (S8), with
$n=m$, is halted.  The right side is therefore halted too.  The reverse
halting-witness transformation is simply $m\mapsto m$; no least-time
search, classical existential principle, or Markov principle is used.
Consequently
\[
 \mathrm{iPRA}\vdash
 \operatorname{Halt}(F(P),C_0)\longleftrightarrow
 \operatorname{Halt}(P,C_0),
\]
and intuitionistic logic alone gives the equivalence of their nonhalting
sentences.  For first halting times, when they exist, (S4)--(S5) also show
that no halt is concealed inside a skipped segment.

\subsection{Compatibility with the B history selector}

PA194ORHB differs from PA194ORH in exactly four native rows:
\[
\begin{array}{c|c|c}
\text{row}&\text{ORH}&\text{ORHB}\\\hline
(87,0)&0R301&0R88\\
(88,0)&0R122&0L301\\
(301,0)&0R45&0R122\\
(301,1)&1R103&1R301.
\end{array}                                                    \tag{S9}
\]
None is a row used in (S4)--(S5), and none changes an incoming edge into
11, 26 or 101.  Therefore the entire proof (S1)--(S8) applies to ORHB
without an additional history assumption.  The B transformation has its
own separately checked relative simulation; composing that theorem with
(S8) is legitimate.  In particular, the F proof covers B's interior
configurations as well as its simulation boundaries.

For blank input we obtain both relative iPRA results
\[
 \neg\operatorname{Halt}(\mathrm{ORHF})
       \longleftrightarrow\neg\operatorname{Halt}(\mathrm{ORH}),
\]
\[
 \neg\operatorname{Halt}(\mathrm{ORHBF})
       \longleftrightarrow\neg\operatorname{Halt}(\mathrm{ORHB}).
\]
Composing with the B relative theorem, the ORH relative theorem, and the
explicitly inherited revised AX11H written calibration yields the same
$iPRA$ equivalence with the fixed PA-consistency predicate.  This uses
ordinary mathematical proofs and finite certificates; no new completed
end-to-end Lean or Rocq theorem is claimed.

\subsection{The independent finite certificate}

The accompanying \texttt{verify\_shortcuts.py} reads both complete parents,
checks all incoming-edge sets in (S2), checks the exact B difference (S9),
and constructs the two F children by exactly the specified rows.  For each
of the four tables it checks preservation of (S1) over every control and
every assignment to the seven local cells $[-3,3]$ satisfying (S1): 24,656
cases per table.  This is exhaustive for the formula because after a
one-cell move every cell it mentions is still in that window.  It also
checks (S4) on all 16 compatible seven-cell assignments and (S5) on all
32 compatible assignments, preserving the entire local tape and never
reading outside the window.  The other 386 rows are identical.

These finite checks establish the candidate-specific closed Boolean
facts.  The induction and clock construction above explain why those
facts prove a uniform iPRA theorem, including for model-internal run
lengths.  The certificate is not used as a simulation sample standing
in for that argument.  Optional child-table arguments make the checker
compare the complete delivered child bytes after parsing, rather than
only reconstructing a proposed child in memory.
